\documentclass[11pt]{extarticle}


\usepackage{tgadventor}
\renewcommand*\familydefault{\sfdefault} %% Only if the base font of the document is to be sans serif
\usepackage[T1]{fontenc}
\usepackage[margin=3.5cm]{geometry}
\usepackage{array}
\usepackage{draftwatermark}
\SetWatermarkText{Mock specimen}
\SetWatermarkScale{.5}

\setlength{\parindent}{0em}
\setlength{\parskip}{.5em}
\setlength{\headheight}{55pt}

\usepackage{tabularx}


%Define fake commands for values to be replaced by the pre-processor
\newcommand{\data}[1]{}
\newenvironment{dataiter}[1]{}{}


\usepackage{fancyhdr}
\pagestyle{fancy}
\fancyhf{} 

\fancyhead[L]{\parbox[b]{0.65\textwidth}{%
  \textbf{\LARGE Toronto Cancer Hospital} \\[0.5em]
 \large Division of Tumor \\ 
  Sequencing and \\ 
  Diagnostics
}}
\fancyhead[R]{\parbox[b]{0.3\textwidth}{%
  \raggedleft
  123 Main St. West, Toronto, ON, A4B5G9 \\[0.5em]
  info@och.on.ca \\
  fax: 416-456-7890\\
  phone: 437-416-6470
}}

\fancyfoot[L]{\textbf{MRN: ~~~~~~~~~~~ Date: 2023-12-08 01:14:00}}
\fancyfoot[R]{\textbf{Page \thepage\ }}

\renewcommand{\footrulewidth}{0.25pt}

\begin{document}
\hfill

\hrule
\begin{center}
{\Huge \bf Draft {\uppercase Molecular Genetics LABORATORY RESULTS}
\end{center}
\hrule

\begin{tabular}{p{9cm} p{6cm}}
\parbox[t]{9cm}{
  \textbf{\Large Patient Info:} \\[0.5em]
  Name: \\
  DOB: \\
  Sex: \\
}
&
\parbox[t]{6cm}{
  \textbf{\Large Patient Information:} \\[0.5em]
  Physician Name:\\
  Health Card:\\
  MRN \#: \\
  Clinic: Peterborough Regional Health Centre (Peterborough) %Lei struggle area
}
\end{tabular}
\vspace{2em}
\hrule
\newcolumntype{C}{>{\centering\arraybackslash}X}
{\bf
\begin{tabularx}{\textwidth}{C C C}
Procedure Date & Accession Date & Report Date \\
2023-12-08 01:14:00 & 2023-12-09 01:26:05 & 2023-12-10 00:29:33
\end{tabularx}
}
\vspace{1.5em}

{\bf \Large Report Details}
\\ \\
Genome Reference - NCBI Build 36.1
Sequencing Range - Whole transcriptome sequencing (WTS), Fusion analysis\\ \\
Referral Reason - Clinical\\ \\
Analysis Location - Kingston Health Sciences Centre (Kingston)


{\bf \Huge Table of findings:}
\vspace{1em}
\newcolumntype{W}{>{\hsize=0.85\hsize}X} %mid column
\newcolumntype{Y}{>{\hsize=1.75\hsize}X} % Wider column
\newcolumntype{Z}{>{\hsize=0.35\hsize}X} % Narrower column

\begin{tabularx}{\textwidth}{Z|Y|X}
{\bf \large Gene} & {\bf \large Information} & {\bf \large Interpretation} \\
\hline

NACA & 3 c.5405T>C p.Leu1802Pro homozygous & Likely pathogenic \\ 
 
%BRCA1 & 11 c.123A>T p.Lys41* Heterozygous & Pathogenic \\  %test compile
&& \\
%VHL & 5 c.499c>T p.Arg167Trp heterozygous & Uncertain clinical significance %also test
\end{tabularx}
\vspace{3em}
\newcolumntype{S}{>{\hsize=0.15\hsize}X}

\begin{tabularx}{\textwidth}{S X}
{\bf \large Summary: } & One likely pathogenic variant detected.. The details regarding the specific mutations are included below. 
\end{tabularx}

\newpage
\vspace{2em}
{\Huge \bf Variant Interpretation: } 
\newline
\newline The interpretation of these variants is as follows: One likely pathogenic variant   was detected in the sample. 

 \vspace{2em}{\bf \large Variant 1 of 1 }
      \newline \vspace{2em} \begin{tabularx}{\textwidth}{C C C C} 
 &&&\\ Gene & Variant & Amino & Zygosity\\ NACA  &  c.5405T>C & p.Leu1802Pro  &  homozygous 
\end{tabularx} \vspace{2em} {\bf Likely pathogenic } \newline
The g.56717709T>C variant occurs in chromosome chr12 , within the NACA gene, and it causes g.56717709T>C change at position 1802 in exon 3 , forming p.Leu1802Pro . This mutation has been identified in 50 families. It has a population frequency of 2.34e-04 (105 alleles in 448325 total alleles tested), indicating it is a rare variant in the general population.
causes an amino acid substitution, which replaces leucine with proline .
ClinVar and other genomic databases report the NACA c.5405T>C variant as clinically relevant based on aggregated evidence. 

 This variant is classified as likely pathogenic. It is believed to negatively impact protein function and may 
        play a role in disease development in affected individuals. 

The affected nucleotide lies within a region 
  that is highly conserved across vertebrate species, which suggests functional importance and evolutionary constraint. 


      This variant is implicated in oncogenesis and other disease processes according to ClinVar records  (VCV accession: VCV004163031). 

Supporting studies and case reports can be found 
  in the scientific literature. Relevant PubMed references include: 689373919, 803116645, 266468244, 451944838, 559393418, 540697506, 781355588 . \newpage
In accordance with existing evidence, this variant is therefore classified as a likely pathogenic variant.


{\huge \bf \underline{Test Details:}} \newline
\newline
{\bf \Large Genes Analyzed: } RARA, total. MARK1, total. VTI1A, total. VHL, total. DACH2, total. NACA, total. MXRA5, total. ERG, total. CNTRL, total.  
9 total genes tested. 

{\bf \large mRNA ref sequences tested NM_...:}
NM\_000964.4 NM\_018650.5 NM\_145206.4 NM\_000551.4 NM\_053281.3 NM\_001365896.1 NM\_015419.4 NM\_182918.4 NM\_007018.6 

\newline \newline
{\Large \bf Recommendations \newline}
A precision oncology approach is advised. These mutations are known oncogenic drivers, linked to constitutive pathway activation. Targeted therapies, including hormone-correcting agents, may be considered, guided by clinical judgment. PI3K inhibitors could be explored in trials for PIK3CA-mutated cases. Germline testing is not indicated, as all mutations are consistent with somatic events. A multidisciplinary tumour board review is recommended to integrate findings into care. Additional testing may be pursued at the physician’s discretion. 
\newline 
\newline
{\Large \bf Methodology \newline}
Ribo-Zero RNA was sequenced with Whole transcriptome sequencing (WTS) and was analyzed using Fusion analysis covering all coding exons and adjacent intronic regions. Target enrichment was performed with hybrid capture (Twist Bioscience), followed by Illumina NextSeq sequencing. Reads were aligned to GRCh37 using BWA-MEM, and variants called with GATK. Annotation was performed in VarSeq using population databases, predictive algorithms, and ClinVar. CNVs were assessed with CNVkit and confirmed by MLPA when applicable. Regions with pseudogene interference, such as PMS2, were validated using long-range PCR and Sanger sequencing. Mean read depth exceeded 300x, with a minimum threshold of 50x. Analytical sensitivity is >99\% for SNVs/indels and >95\% for exon-level CNVs. {\bf Only variants classified as pathogenic, likely pathogenic, or variants of uncertain significance (VUS) are reported}, per ACMG/AMP guidelines (PMID: 25741868).

{\Large \bf Limitations \newline}
This test was developed and validated in a certified clinical laboratory. Limitations include reduced sensitivity in pseudogene regions (e.g., PMS2, CHEK2), and inability to detect certain structural variants (e.g., MSH2 inversions), deep intronic changes, or low-level mosaicism. PMS2 exons 11–15 are not analyzed. Interpretations reflect current knowledge and may be updated as new evidence emerges.

\newpage
{\huge Report Electronically Verified and Signed by: }

\it{Missing} maf
\end{document}